% Labels of the draft are printed as \dref{label} (texttt), since this note is compiled on its own.
\providecommand{\dref}[1]{\texttt{\detokenize{#1}}}
\def\v#1{\mathbf{#1}}% vectors as in the draft
\section{Referee 3, Major 7 and the items on Sections 4.3 and 7.3: the cubic terms on exact charge states}\label{sec:ref3cubic}
% Scripts: notes/scripts/ref3_cubic/ (colour.py, quark_dipole.py, singlequark_table.py, magnitudes.py, emission.py,
% cparity.py, ls_check.py; logs *.log). Run with python3 -I from that directory.
% Sources read: inputs/reviews/Referee_report_3.txt; draft sec_setup.tex (lcpt, A3threerho), sec_observable.tex (Groups, rows),
% sec_rapidity.tex (4.2, 4.3, 4.5), sec_finite.tex (fin_singlet incl. A2d, cubicavg, (2rho)/(3rho) rows, Deltas, sodd),
% app_finite.tex (app:fin_sym, rows II.I, III.II), sec_checks.tex (odd terms, 7.3); notes/fix_singlet.tex, fix_dabc.tex,
% ref3_prescriptions.tex (sibling note, Major 5); Li--Skokov arXiv:2104.01879 and 2102.01594, McLerran--Skokov
% arXiv:1611.09870, Kovchegov--Skokov arXiv:1802.08166 (sources from export.arxiv.org).

\noindent\textbf{Result in one paragraph.}
(a) For a single quark and for a colour-singlet $q\bar q$ dipole every cubic colour structure of the finite terms was evaluated exactly in charge space (explicit $SU(N_c)$ matrices, ordered products as written in Appendix~\dref{app:fin_sym}, $N_c=2,3,4$, all charge assignments): the exact expectation value equals symmetrised part $+$ contact terms of Eq.~(\dref{eq:cubicavg}) to $5\times10^{-14}$ (largest relative deviation over all cases; $3\times10^{-16}$ for $SU(4)$). For a single quark the whole content is $\frac1{N_c}{\rm tr}(t^at^bt^c)=\frac{1}{4N_c}(d^{abc}+if^{abc})$: the symmetrised part is the $d^{abc}$ half, the contact part the $if^{abc}$ half; $o=\frac{(N_c^2-4)(N_c^2-1)}{4N_c^2}=C_F\frac{N_c^2-4}{2N_c}$ ($=\frac{10}{9}$ for $SU(3)$), $n=C_F$. For the dipole $o=o_q\,s_is_js_k$ and $n=C_F\,s_is_j$ ($s=+1$ at the quark, $-1$ at the antiquark): the $k$-odd part is the part odd under the exchange $q\leftrightarrow\bar q$ (the orientation of the dipole), the $k$-even cubic part is even under it. Both parts are of the same order as the quadratic terms; there is no density enhancement of either for a quark or a dipole (colour weights: symmetrised/LO $\approx0.12$, contact/LO $\approx0.8$ for $SU(3)$ and a dense random target). (b) The ordering is not a choice: it is the order in which the coefficients of $\bar\Omega^\dagger\Omega$ and of the conjugate amplitude act, fixed by the light-cone time ordering (Eq.~\dref{eq:lcpt}, e.g.\ $[\mf A^{(3)}]_{\rho\rho\rho}=\mf A\mf N_2$). The decomposition is an operator identity; the symmetrised part is the \emph{same operator for every ordering} of the three charges, so all ordering dependence sits in the contact terms; and under (A2) the two parts are the C-odd and C-even parts of the projectile, $\frac12[\sigma(P)\mp\sigma(P^{\rm C})]$, a characterisation that does not refer to an ordering (verified on the quark/antiquark and dipole/conjugate dipole pairs to $8\times10^{-15}$). (c) Li--Skokov's first saturation correction $n^{(4)}\propto\rho_P^3$ (their $M_{(1)}^*M_{(3)}+{\rm c.c.}$) is the classical counterpart of the symmetrised cubic terms. Agreement: no contribution for a Gaussian projectile, non-zero for a C-odd three-charge correlator; their two-Wilson-line term has zero $d^{abc}$ projection, as the corresponding structures here; the odd part needs three Wilson lines (C-odd target, $\mf G$); in the paper the four-line invariants $\mf Q_d$ (and $\mf Q$ in the contact terms) cancel between Rows~II.I and III.II (new observation, $10^{-15}$), so the odd part has at most three Wilson lines, as in Li--Skokov and McLerran--Skokov. \emph{Disagreement:} Li--Skokov's $n^{(4)}$ has a $k$-even part configuration by configuration (their $\mathcal E_1$ term and ${\rm Re}\,\mathcal E_2$), and the $d^{abc}$ projection of ${\rm Re}\,\mathcal E_2$ is not removed by any symmetry (antisymmetrised kernel non-zero at all of 4000 random points, median $0.54|\mathcal E_2|$); in the draft the symmetrised cubic part is purely $k$-odd. This is consistent with the finding of the sibling note \texttt{ref3\_prescriptions} that the principal-value dressing of $\mf B_3$ drops an in-phase term $\Delta N$ of Row~I.III whose symmetrised part is $k$-even; a kernel-level comparison was not possible (no momentum-space form of the finite terms). (d) The draft (all versions in the repository, source and PDF) already has $\frac{g^2}{4\pi^3}C_FK_D=\frac{\alpha_sC_F}{\pi^2}K_D$; the referee's $\frac{g^2}{4\pi^4}$ is not in the text. We re-derived it from the paper's own Eq.~(\dref{eq:N2}) and verified Eq.~(7.2) exactly at finite $N_c$ (sympy), with ``LO prefactor'' $=\frac{\alpha_s}{4\pi^4k^+}$ (the likely source of the confusion); with $\frac{g^2}{4\pi^4}$ the result would be off by $1/\pi$. (e) Drop-in text for Section~4.3 with the C-parity argument as the primary reason; the precise statement is that C-invariance forbids $o$ in the \emph{projectile} evolution coefficient and in $H_{\rm JIMWLK}\mathcal O_{\rm LO}$; C alone does not forbid $o\times$(C-odd target), which is exactly what the finite terms contain, so the factorised form of the logarithms is also needed; the numerical cancellation confirms it and is stronger (it also removes the C-even $f^{abc}g_f$ part).

% ------------------------------------------------------------------------------------------
\subsection{(b) Where the ordering comes from, and why the split is unique}\label{r3c:b}
\paragraph{Origin of the ordering.} Every row is a term of $\langle\mathcal A^\dagger\mathcal A\rangle$, Eqs.~(\dref{eq:N_amplitudes}), (\dref{eq:Groups}), with the conjugate amplitude to the left; inside each amplitude the coefficients of $\bar\Omega^\dagger\Omega\ket0$ stand in the order in which they act ($\Omega$ first, rightmost), and inside each coefficient the matrix elements stand in the order of the light-cone time ordering, ``the vertex acting later stands to the left'' (Eq.~\dref{eq:lcpt}); the norm factor stands to the right of the emission, $[\mf A^{(3)}]_{\rho\rho\rho}=\mf A\,\mf N_2$ (Eq.~\dref{eq:A3threerho}). Examples: Row~I.III $=\mathrm{LO}^\dagger\times\bar{\mf B}_3^\dagger\mf A$ gives $\rho(x')\rho(x)\rho(y)$ [conjugate LO, then the rotated charge of $\bar{\mf B}_3^\dagger$, then the charge of $\mf A$]; Row~I.II $N_{2b}$, I.V, III.III have one $\{\rho,\rho\}$ (the $\{\rho,\rho\}$ part of $\mf B_2$ or the two-$\rho$ part of $\mf A^{(3)}$). With the wrong order of $\mf A$ and $\mf N_2$ the draft notes a spurious $\ln(\vee/\Lambda)$ two-$\rho$ term (Eq.~\dref{eq:A3threerho} and below). So the ordered product of each row, and hence its expectation value in any projectile state, is unique.
\paragraph{The decomposition is an identity, and its symmetrised part is ordering-independent.} Eq.~(\dref{eq:symmetrise}) is an operator identity. Its first term $C^{abc}\{\rho^a(p_1)\rho^b(p_2)\rho^c(p_3)\}_{\rm sym}$ is literally the same operator for all six orderings of the three factors $C^{abc}\rho^a(p_1)$, $\rho^b(p_2)$, $\rho^c(p_3)$; only the commutator (contact) terms change with the ordering. A different (wrong) ordering would therefore change the contact terms and leave the symmetrised part, and with it the $k$-odd part of the spectrum, unchanged. Verified on the dipole (\texttt{quark\_dipole.py}, Row~I.III$(z)$, assignment $(q,q,\bar q)$): the six orderings give $0.0957\pm0.3784\,i$, with the same symmetrised part $0.0957$ (spread $3\times10^{-17}$) and contact parts $\pm0.3784\,i$; each equals sym $+$ contact to $4\times10^{-16}$.
\paragraph{Physical (ordering-free) characterisation.} Under charge conjugation of the projectile, $\rho^at^a\to-(\rho^at^a)^T$, the symmetrised three-point function goes into $-\eta_a\eta_b\eta_cS^{abc}$, i.e.\ $d^{abc}g_d\to-d^{abc}g_d$ and $f^{abc}g_f\to+f^{abc}g_f$, while the commutator terms ($\propto f^{abc}n$) are unchanged ($n$ is C-even). Hence, under (A1), (A2),
\begin{align}\label{r3c:Csplit}
\text{symmetrised part}=\tfrac12\big[\langle C\rho\rho\rho\rangle_P-\langle C\rho\rho\rho\rangle_{P^{\rm C}}\big],\qquad
\text{contact part}=\tfrac12\big[\langle C\rho\rho\rho\rangle_P+\langle C\rho\rho\rho\rangle_{P^{\rm C}}\big],
\end{align}
and, by the C-invariance of the whole result [$\langle C\rho\rho\rho\rangle_{P^{\rm C}}[V^*]=\langle C\rho\rho\rho\rangle_P[V]$, verified exactly for all 13 structures on the dipole, \texttt{cparity.py}], the same split is the target-C-odd/even split, $\frac12[\sigma(V)\mp\sigma(V^*)]$, and, by the kernel parity of Section~\dref{sec:fin_singlet}, the $k$-odd/even split. Three definitions that do not mention an ordering coincide with the split ``symmetrised $+$ contact''. Any other reference ordering would not separate C-odd from C-even pieces (e.g.\ for a quark $\langle t^at^bt^c\rangle=\frac1{4N_c}(d^{abc}+if^{abc})$ mixes them). [If $g_f\neq0$, which (A2) excludes, the $f^{abc}g_f$ part is symmetrised and $k$-odd but C-even; then the $k$-odd and C-odd splits differ by it.]

% ------------------------------------------------------------------------------------------
\subsection{(a) Single quark, exactly}\label{r3c:quark}
With the normalisation of Eq.~(\dref{eq:rho}) a quark at $x_0$ is $\rho^a(x)=t^a\delta^{(2)}(x-x_0)$ [the referee's $g\,t^a$ corresponds to keeping $g$ in $\rho$; then $n\to g^2n$, $o\to g^3o$], with the colour-averaged state $\varrho=\mathbb 1/N_c$ (globally colour invariant). Exactly,
\begin{align}\label{r3c:q3}
\langle\rho^a(p_1)\rho^b(p_2)\rho^c(p_3)\rangle=\tfrac1{N_c}{\rm tr}(t^at^bt^c)\prod_i\delta^{(2)}(p_i-x_0)=\frac{d^{abc}+if^{abc}}{4N_c}\prod_i\delta^{(2)}(p_i-x_0)\,,
\end{align}
so $S^{abc}=\frac{d^{abc}}{4N_c}\prod\delta$, $g_f=0$ (A2 holds), and
\begin{align}\label{r3c:oq}
o=o_q\prod_i\delta^{(2)}(p_i-x_0)\,,\quad o_q=\frac{d^{abc}d^{abc}}{4N_c}=\frac{(N_c^2-4)(N_c^2-1)}{4N_c^2}=C_F\,\frac{N_c^2-4}{2N_c}\,,\qquad n=C_F\,\delta\delta\,,
\end{align}
[$d^{abc}t^at^bt^c=C_F\frac{N_c^2-4}{2N_c}\mathbb 1$, $f^{abc}t^at^bt^c=\frac{i}{2}N_cC_F\mathbb 1$; $o_q=\frac{10}{9}$, $2.8125$ for $N_c=3,4$; $0$ for $SU(2)$]. All three charge points sit at $x_0$, so all Wilson lines at charge points are $U(x_0)$ and the kernels are taken at the triple coincidence $x=x'=y=x_0$ ($y'$, $z$, $w$, $w'$ remain). For each type:
\begin{align}\label{r3c:qtypes}
\text{ord: }\ \langle C\rho\rho\rho\rangle=\frac{\Delta_{\rm loc}+i\Phi_{\rm loc}}{4N_c}\,,\qquad
\text{A23, A12: }\ \langle C\rho\rho\rho\rangle=\frac{\Delta_{\rm loc}}{4N_c}\quad\big[{\rm tr}(t^a\{t^b,t^c\})=\tfrac12d^{abc}\big]\,,
\end{align}
with $\Delta=\sum_{\rm orderings}C^{abc}d^{abc}$ and $\Phi=C^{abc}f^{abc}$ at the coincidence. The split of Eq.~(\dref{eq:cubicavg}) reproduces this term by term: symmetrised part $\frac{N_c}{(N_c^2-4)(N_c^2-1)}\Delta\,o_q=\frac{\Delta}{4N_c}$; contact part $\frac{i\Phi}{2(N_c^2-1)}C_F(1+1-1)=\frac{i\Phi}{4N_c}$ (ord) and $\frac{i\Phi}{N_c^2-1}C_F(1-1)=0$ (A23, A12). For the antiquark ${\rm tr}(\bar t^a\bar t^b\bar t^c)/N_c=\frac{-d^{abc}+if^{abc}}{4N_c}$: the symmetrised part changes sign, the contact part does not, Eq.~(\ref{r3c:Csplit}). Closed forms at the coincidence ($a\equiv x_0$; $\mf D,\mf F,\mf G,\mf Q_d$ of the draft; $\mf F(a,a,c)=N_c\mf D(a,c)$, $\mf Q(a,b;c,a)=\mf F(a,b,c)$ used), verified against the exact trace (\texttt{singlequark\_table.py}, $SU(3)$, $SU(4)$):
\begin{center}\small\renewcommand{\arraystretch}{1.25}
\begin{tabular}{|l|l|l|l|}\hline
structure & type & $\Delta_{\rm loc}$ (symmetrised, $k$-odd) & $\Phi_{\rm loc}$ (contact, $k$-even; ord only)\\\hline
I.II $N_{2b}$ & A23 & $0$ & ---\\
I.III $(u{=}z,w)$ & ord & $\mf G(a,u,w')$ & $N_c(N_c^2{-}1)-N_c\mf D(a,w')-N_c\mf D(a,u)+\mf F(a,u,w')$\\
I.V $N_1$ & A23 & $-2[\mf G(x,a,z)+\mf G(x,z,w')]$ & ---\\
II.I $N_{1a}$ & A12 & $2\mf G(a,z,w')-2\mf Q_d(z,a;a,w)$ & ---\\
II.I $N_{1b}$ & A23 & $2\mf G(a,z,w)-2\mf Q_d(z,a;a,w')$ & ---\\
III.II & ord & $\mf Q_d(z,a;a,w)$ & $N_c(N_c^2{-}1)+N_c\mf D(a,w)-N_c\mf D(a,z)-\mf F(a,z,w)$\\
III.III$_1$ $(y')$, $(z)$ & A23 & $2\mf G(y',z,w)-2\mf G(a,z,w)$, $\ -2\mf G(a,z,w)$ & ---\\
III.V $(v{=}y',z)$ & ord & $\mf G(a,v,w)$ & $-N_c(N_c^2{-}1)+N_c\mf D(a,v)+N_c\mf D(a,w)-\mf F(a,v,w)$\\
III.VI $(y')$ & ord & $-\mf G(a,y',z)$ & $N_c(N_c^2{-}1)-N_c\mf D(a,y')-N_c\mf D(a,z)+\mf F(a,y',z)$\\
III.VI $(z)$ & ord & $0$ & $N_c(N_c^2{-}1)-N_c\mf D(a,z)$\\\hline
\end{tabular}\end{center}
(In Row~I.V the point $x$ is not a charge point.) \textbf{Physical content for a single quark.} (i) The symmetrised part is the C-odd part, $\frac12[\sigma(q)-\sigma(\bar q)]$, and needs a C-odd target ($\mf G$, $\mf Q_d$; $O(\epsilon^3)$ for a weak target); the contact part is $\frac12[\sigma(q)+\sigma(\bar q)]$. (ii) Neither is enhanced: $\rho\sim t^a$, so $\rho^3$ and $\rho^2$ terms are of the same size; with the quadratic weight $\langle C^{ab}\rho^a\rho^b\rangle=C^{aa}\frac{C_F}{N_c^2-1}=\frac{C^{aa}}{2N_c}$ the cubic weights carry $\frac{1}{4N_c}$ with colour factors of the same type ($\mf D,\mf F$ in $\Phi$). For a quark the cubic finite terms are simply part of the $O(\alpha_s)$ correction, of the same order as the quadratic finite terms; the language ``first saturation correction'' does not apply. (iii) All anticommutator rows (I.II $N_{2b}$, I.V, II.I, III.III) have no contact term for a single quark; Row~I.II $N_{2b}$ vanishes identically. (iv) In the sum of Rows~II.I and III.II the $\mf Q_d$ cancel (Sec.~\ref{r3c:c}): $\frac12\Delta_{\rm loc}^{1a}+\Delta_{\rm loc}^{\rm III.II}=\mf G(a,z,w')$. (v) The kernels are needed at the triple coincidence $x=x'=y=x_0$; their integrability there was not examined (it is the open point (iii) of the draft's Remarks for a local $o$).

% ------------------------------------------------------------------------------------------
\subsection{(a) Colour-singlet dipole, exactly}\label{r3c:dipole}
Quark at $x_0$, antiquark at $x_1$, $\rho^a(x)=Q_0^a\delta^{(2)}(x-x_0)+Q_1^a\delta^{(2)}(x-x_1)$ with $Q_0=t\otimes\mathbb 1$, $Q_1=\mathbb 1\otimes(-t^T)$ and the singlet $\ket s=\sum_i\ket{ii}/\sqrt{N_c}$, $(Q_0+Q_1)\ket s=0$. Exactly (\texttt{quark\_dipole.py}),
\begin{align}\label{r3c:dip}
n(p,q)=C_F\sum_{i,j}s_is_j\,\delta_{p,x_i}\delta_{q,x_j}\,,\qquad o(p_1,p_2,p_3)=o_q\sum_{i,j,k}s_is_js_k\,\delta_{p_1,x_i}\delta_{p_2,x_j}\delta_{p_3,x_k}\,,\qquad s_0=+1,\ s_1=-1\,,
\end{align}
($\delta_{p,x}\equiv\delta^{(2)}(p-x)$), $g_f=0$ to $10^{-16}$; the dipole has $o\neq0$: it is not C-symmetric, its C-conjugate is the dipole with $q$ and $\bar q$ exchanged ($o\to-o$, $n\to n$). For a cubic structure with charge points at $(x_i,x_j,x_k)$,
\begin{align}\label{r3c:dipsplit}
\langle C\rho\rho\rho\rangle_{ijk}=\frac{1}{4N_c}\Big[s_is_js_k\,\Delta_{ijk}+i\,\Phi_{ijk}\,c_{ijk}\Big]\,,\qquad
c_{ijk}=\begin{cases}\delta_{ij}s_is_k+\delta_{jk}s_is_j-\delta_{ik}s_is_j&\text{ord}\\ 2(\delta_{ij}s_is_k-\delta_{ik}s_is_j)&\text{A23}\\ 2(\delta_{jk}s_is_j-\delta_{ik}s_is_j)&\text{A12}\end{cases}
\end{align}
($\delta_{ij}$: Kronecker delta of the sites). Exact evaluation of all 13 structures for all 8 assignments, random Wilson lines (equal at coinciding sites): exact $=$ Eq.~(\ref{r3c:dipsplit}) to $3\times10^{-14}$ ($SU(3)$; $3\times10^{-16}$ for $SU(4)$; for $SU(2)$ the symmetrised part is zero and the contact part alone reproduces the exact value, to $5\times10^{-14}$). \textbf{Which parts are odd, which even.} The symmetrised part carries $s_is_js_k$: it changes sign under $q\leftrightarrow\bar q$ (orientation of the dipole) and under $V\to V^*$, and it is the $k$-odd part. The contact part carries products of two $s$'s: it is even under $q\leftrightarrow\bar q$ and $V\to V^*$, and $k$-even. In words: for a dipole the $k$-odd cubic term is the part of the spectrum that distinguishes the quark end from the antiquark end, $\frac12[\sigma(x_0,x_1)-\sigma(x_1,x_0)]$; it vanishes for an orientation-averaged dipole ($o=0$), and needs a C-odd target. The $k$-even cubic term survives the orientation average. \textbf{Magnitude} (colour weights only, without kernels; \texttt{magnitudes.py}; r.m.s.\ over random targets and over the assignments, in units of the r.m.s.\ LO weight $\langle[U(y')-U(x')]^{ce'}[U(y)-U(x)]^{ce}\rho^{e'}(x')\rho^e(x)\rangle$):
\begin{center}\small\renewcommand{\arraystretch}{1.2}
\begin{tabular}{|l|c|c|c|c|}\hline
dense (Haar-random) target & $SU(3)$ & $SU(4)$ & $SU(5)$ & $SU(6)$\\\hline
symmetrised ($k$-odd) / LO, mean over 13 structures (range) & $0.12$ ($0.05$--$0.23$) & $0.08$ ($0.03$--$0.15$) & $0.05$ ($0.02$--$0.09$) & $0.04$ ($0.01$--$0.07$)\\
contact ($k$-even) / LO, mean (range) & $0.76$ ($0.38$--$1.07$) & $1.05$ ($0.51$--$1.45$) & $1.12$ ($0.20$--$1.87$) & $1.26$ ($0.15$--$2.19$)\\\hline
\end{tabular}\end{center}
(6 target configurations for $SU(3)$, $SU(4)$, 3 for $SU(5)$, $SU(6)$; statistical accuracy about $20\%$.) The contact weights of the ordered structures grow roughly like $N_c/3$, those of the anticommutator structures do not. For a weak target ($V=e^{i\epsilon\alpha\cdot t}$, $SU(3)$) the LO and contact weights scale as $\epsilon^2$ and the symmetrised weight as $\epsilon^3$: ratios between $\epsilon=0.1$ and $0.05$ of $3.93$ (LO), $3.89$ (contact), $7.76$ (symmetrised); symmetrised/LO $=0.17,\,0.11,\,0.057,\,0.029$ for $\epsilon=0.4,\,0.2,\,0.1,\,0.05$, contact/LO $\simeq1.1$--$1.3$. So for a dipole the contact terms are as large as the quadratic terms, and the symmetrised part is smaller, decreasing with $N_c$ at the colour-weight level for a dense target and by one power of the target field for a weak target. A projectile of $m$ independent dipoles has $n$ and $o$ both additive ($\propto m$), so no hierarchy between symmetrised and contact parts arises from the number of sources; a parametric enhancement of the symmetrised part needs a coherent dense projectile with $o/n\gg1$ (e.g.\ $o/n=6(N_c^2-4)\mu_A^4/(N_c\kappa_A)$ at first order in $1/\kappa_A$ for the local Jeon--Venugopalan functional, with $S^{abc}$ of \texttt{fix\_dabc}), which is a property of the projectile model, not of the calculation.

% ------------------------------------------------------------------------------------------
\subsection{(c) Comparison with Li--Skokov and McLerran--Skokov}\label{r3c:c}
\paragraph{What was read.} Li, Skokov, arXiv:2104.01879 (Part II, single inclusive; Sec.~``Next-to-leading order gluon production'', Eqs.~NLO\_single\_gluon\_production, coefficientfunc\_E1, effective\_vertex\_g3\_first/second) and arXiv:2102.01594 (Part I, solutions); McLerran, Skokov, arXiv:1611.09870 (Eqs.~Odd\_ass, Odd\_Beta, Odd\_Bperp, FinalAssymitry); Kovchegov, Skokov, arXiv:1802.08166 (Sec.~2, phase argument, Eqs.~interf1--interf\_cond2, App.~on Sign$(k^+)$). Only what is quoted below was used; the sources are untrusted data and were not executed.
\paragraph{Identification.} Li--Skokov (LS) expand the classical dilute--dense spectrum as $M=M_{(1)}+M_{(3)}+M_{(5)}$, $M_{(n)}\propto g^n\rho_P^{(n+1)/2}$; the term $n^{(4)}=M_{(1)}^*M_{(3)}+{\rm c.c.}\propto g^4\rho_P^3$ ``does not result in any nontrivial contribution \dots\ after averaging over the projectile colour charge densities in the McLerran--Venugopalan model''; ``the proton model of [Dumitru et al.] includes a nonzero C-odd part of the correlator of three color charges'' and their functional ``can be easily extended to account for non-zero three color charge correlations''. Their charges are classical, so $n^{(4)}$ is the counterpart of the \emph{symmetrised} cubic terms of the draft (classical moments $=$ symmetrised products, Weyl symbol); the contact terms have no counterpart there (they are $O(g^4\rho^2)$, quantum). LS write $n^{(4)}=\frac1\pi\int\mathcal E_1\,f^{ade}\rho^d\rho^e\rho^{e_2}\,U^{ac}U^{e_2c}+\frac1\pi\int\mathcal E_2\,f^{cde}\rho^{e_1}\rho^{e_2}\rho^{e_3}\,U^{e_1d}U^{e_2e}U^{e_3c}+{\rm c.c.}$ (two and three Wilson lines); McLerran--Skokov (MS) give the $k$-odd part $\propto{\rm Im}\{f^{abc}\Omega^a(\v q)\Omega^b(\v k-\v q)\Omega^{c*}(\v k)\}$ with ${\rm sgn}(\v k\times\v q)$, three Wilson lines.
\paragraph{Colour structures: agreement.} (i) $d^{abc}$ projection of the two-line term: $f^{ade}d^{dee_2}(UU^T)^{ae_2}=0$ (\texttt{ls\_check.py}, $<10^{-16}$): it does not contribute under (A2). In the draft the same mechanism removes the corresponding structures (two legs of $f$ on charges, e.g.\ $f^{hbc}U^{ah}(w)\{\rho^b,\rho^c\}$ in I.II $N_{2b}$, \texttt{fix\_dabc}). (ii) The three-line term projects onto the $\mf G$-type invariant, totally antisymmetric under the simultaneous exchange of (Wilson-line argument, charge argument) pairs (checked), C-odd in the target, $O(\epsilon^3)$ for a weak target: the same as the $\mf G$ of Eq.~(\dref{eq:Gfund}); MS: ``essentially they are defined by odderon exchanges''; Kovchegov--Skokov: three scatterings in the target. (iii) LS and MS have at most three Wilson lines at $O(g^4\rho^3)$. In the draft, Rows~II.I and III.II contain the four-line $\mf Q_d$; \emph{they cancel in the sum}: both rows have the same kernel $\mathfrak K_{\mathcal T}(w,w')$ (real), III.II $+$ c.c.\ gives $-4\sin(\v k\cdot(\v w-\v w'))\mathfrak K_{\mathcal T}\Delta_{\rm III.II}$, and
\begin{align}\label{r3c:Qdcancel}
\Delta_{\rm II.I}+\Delta_{\rm III.II}=\mf G(x,x',z)+\mf G(x,z,w')\,,
\end{align}
so the II.I and III.II lines of Eq.~(\dref{eq:sodd}) combine into $-\frac{\mathcal C_{\rm f}}{N_c^2-1}\frac{N_c}{N_c^2-4}\int4\sin(\v k\cdot(\v w-\v w'))\,\mathfrak K_{\mathcal T}(w,w')[\mf G(x,x',z)+\mf G(x,z,w')]\,o$; $\mf G(x,y,x')$ cancels as well. The same happens in the contact terms: the $\mf Q(x,z;w,x')$ and $\mf Q(x,w;z,y)$ of Eqs.~(\dref{eq:sIII1}) and (\dref{eq:sIIIII}) cancel, the sums being $N_c\mf D(w,w')-N_c\mf D(x',w)-\mf F(x,z,w')+\mf F(x,x',z)$ at $y\to x$ and $-N_c\mf D(w,w')+\mf F(x,z,w')-N_c\mf D(x,z)+N_c\mf D(x,w)$ at $x'\to x$ (times $2\cos(\v k\cdot(\v w-\v w'))\mathfrak K_{\mathcal T}|$). Checked with the tensors of Appendix~\dref{app:fin_sym}, $SU(3)$, $SU(4)$: $1.5\times10^{-15}$. This is a structural agreement with the classical results and a simplification of the draft's result.
\paragraph{Parity in $\v k$: disagreement.} In the draft every cubic row is $i\times$(real kernel)$\times$(functions with $f(-s)=f(s)^*$), so the symmetrised part is purely $k$-odd, configuration by configuration. LS: $\mathcal E_1$ is real and even under the reflection of all momenta, hence its contribution with the c.c.\ is $k$-even; $\mathcal E_2={\rm Re}+i\,{\rm sgn}(\v k\times\v p)(\dots)$, both parts reflection-even, so ${\rm Re}\,\mathcal E_2$ gives a $k$-even and ${\rm Im}\,\mathcal E_2$ a $k$-odd contribution; LS state explicitly that the imaginary part is ``responsible for the asymmetry'' and that ``the even part under $\v k\leftrightarrow-\v k$ is computed here for the first time''. Checks (\texttt{ls\_check.py}): $\mathcal E_2$ as printed equals $\Gamma_2^jL^j$ as printed ($10^{-14}$, transcription check); reflection properties exact. After the projection (A2) the only domain symmetry that acts on the three-line term is $(\v p\to\v k-\v p,\ \v p_1\leftrightarrow\v p_2)$, under which the colour factor is odd; the surviving kernel $\mathcal E_2^A=\frac12[\mathcal E_2(\v p,\v p_1,\v p_2,\dots)-\mathcal E_2(\v k-\v p,\v p_2,\v p_1,\dots)]$ has $|{\rm Re}\,\mathcal E_2^A|$ of median $0.54\,|\mathcal E_2|$ at 4000 random points and vanishes at none. Hence LS predict, for $o\neq0$, a $k$-even term $\propto o\times$(three-line C-odd target structure), which the draft does not have. Kovchegov--Skokov's phase argument says the same in LCPT language: the $k$-odd part comes from the part of $M_3$ with a phase difference to $M_1$, the in-phase part gives an even contribution. \emph{Possible origin:} the sibling note \texttt{ref3\_prescriptions} finds that the draft's principal-value dressing of $\mf B_3$ (Row~I.III) drops a term $\Delta N$ of the adiabatic out-state dressing, with a \emph{real} kernel, whose symmetrised three-$\rho$ part is $k$-even (and its commutator part $k$-odd). A term of this kind is what the LS comparison requires. Whether $\Delta N$ reproduces ${\rm Re}\,\mathcal E_2$ (and the event-by-event $\mathcal E_1$) was not checked.
\paragraph{Kernels, dilute target, single quark.} A kernel-level comparison needs the cubic finite terms in momentum space, which the draft does not provide (e.g.\ ``the coordinate-space form of this kernel has not been derived'' for I.II $N_{2b}$, and conversely for the other rows); it was not done. For a single quark (Weyl-symmetrised charges) the colour normalisation agrees by construction: both results are linear in $\langle\{\rho\rho\rho\}_{\rm sym}\rangle=\frac{d^{abc}}{4N_c}$, so the odd part is $\frac{1}{4N_c}\times$(three-line $\mf G$-type factor) $\times$ kernel in both, and the two-line terms drop. A classical calculation is, however, not applicable to a single quark: there the contact part is as large as the symmetrised part (Sec.~\ref{r3c:quark}). In the dilute-target limit both odd parts start at $O(\epsilon^3)$ (checked for all 13 $\Delta$ in \texttt{fix\_dabc}; LS/MS: three-line C-odd structure); the coefficients were not compared.

% ------------------------------------------------------------------------------------------
\subsection{(d) Section 7.3: normalisation of the dipole emission probability}\label{r3c:d}
\paragraph{What the draft says.} The source of \texttt{sec\_checks.tex} (all eight versions in the git history) and the compiled \texttt{main.pdf} have ``the probability to emit the additional gluon at $\v z$ is $\frac{g^2}{4\pi^3}\ln\frac{\vee}{k^+}C_F\,K_D(\v x_1,\v x_2,\v z)$'', i.e.\ $\frac{\alpha_sC_F}{\pi^2}K_D$, which is what the referee asks for. The form $\frac{g^2}{4\pi^4}$ quoted by the referee does not occur; it may come from the factor ``(LO prefactor)'' of Eq.~(7.2), which is $\frac{\alpha_s}{4\pi^4k^+}$ and is not defined in the text.
\paragraph{Independent derivation} (\texttt{emission.py}). (i) From the draft itself: $\mf N_2=-\frac{g^2}{8\pi^3}\ln\frac\vee\Lambda\int_{x,y,z}K(x,y,z)\rho^a(x)\rho^a(y)$ is minus one half of the emission probability, equivalently $\langle\Psi^{\rm LO}|\Psi^{\rm LO}\rangle$ with $\braket{g|g}=\delta\delta$; with $n=C_F[\delta_{11}+\delta_{22}-\delta_{12}-\delta_{21}]$, $\int K\,n=C_FK_D$:
\begin{align}\label{r3c:P}
dP=\frac{g^2C_F}{4\pi^3}\,K_D(\v x_1,\v x_2,\v z)\,d^2z\,\ln\frac{\vee}{k^+}=\frac{\alpha_sC_F}{\pi^2}\,\frac{(\v x_1-\v x_2)^2}{(\v x_1-\v z)^2(\v x_2-\v z)^2}\,d^2z\,\ln\frac{\vee}{k^+}\,,
\end{align}
the standard dipole emission probability; equivalently, Parseval applied to the Gunion--Bertsch density $\frac{\alpha_sC_F}{\pi^2\v k^2}$ checked by the referee. (ii) Charges of the $q\bar qg$ singlet, exactly (explicit matrices, $N_c=2,3,4$): $Q_1^2=Q_2^2=C_F$, $Q_g^2=N_c$, $Q_1\!\cdot\!Q_2=\frac1{2N_c}$, $Q_{1,2}\!\cdot\!Q_g=-\frac{N_c}{2}$ (to $10^{-14}$); $q\bar q$: $Q_1\!\cdot\!Q_2=-C_F$. (iii) Real $-$ virtual (sympy, finite $N_c$): $\sum_{s,s'}Q_s\!\cdot\!Q_{s'}F_{ss'}-C_F(F_{11}+F_{22}-F_{12}-F_{21})=N_c^2[\mathcal D_{1z}+\mathcal D_{z2}-\mathcal D_{12}]$ identically. (iv) Prefactors: with Eq.~(\dref{eq:LOsinglet}) the LO spectrum of the dipole is $\frac{\alpha_s}{4\pi^4(N_c^2-1)k^+}C_F\cdot2N_c\mathcal D_{12}=\frac{\alpha_s}{4\pi^4k^+}\mathcal D_{12}$; then $dP\times\frac{\alpha_s}{4\pi^4(N_c^2-1)k^+}\times N_c^2[\dots]=\frac{\alpha_sN_c}{2\pi^2}\ln\frac\vee{k^+}K_D[\dots]\times\frac{\alpha_s}{4\pi^4k^+}$: ratio to Eq.~(7.2) exactly $1$. With $\frac{g^2}{4\pi^4}$ instead the ratio is $1/\pi$. \textbf{Corrected expression:} the draft's $\frac{g^2}{4\pi^3}C_FK_D$ is correct; only ``(LO prefactor)'' should be defined. Drop-in (D3) below.

% ------------------------------------------------------------------------------------------
\subsection{(e) Section 4.3: C-parity as the reason}\label{r3c:e}
\paragraph{The transformation.} $\mathrm C$: $\rho^at^a\to-(\rho^at^a)^T$, i.e.\ $\rho^a\to-\eta_a\rho^a$ with $(t^a)^T=\eta_at^a$, $\eta_a=\pm1$; $V\to V^*$, i.e.\ $U^{ab}\to\eta_a\eta_bU^{ab}$; soft gluons $A^a\to-\eta_aA^a$. Since $\eta_a\eta_b\eta_cf^{abc}=-f^{abc}$ and $\eta_a\eta_b\eta_cd^{abc}=+d^{abc}$, C preserves Eq.~(\dref{eq:rho}) and every colour-singlet coupling ($\rho^aA^a$, $f^{abc}AAA$, $U^{ab}$ acting on $A^b$, $a^{\dagger a}a^a$): the spectrum satisfies $\sigma[\varrho_P^{\rm C};V^*](\v k)=\sigma[\varrho_P;V](\v k)$. Under C: $n\to n$, $o\to-o$, $g_f\to g_f$; $\mf D,\mf F,\mf Q\to+$, $\mf G,\mf Q_d\to-$ (\texttt{cparity.py}: exact for $SU(3)$, $SU(4)$; on the dipole the full cubic expectation values are invariant under joint C, the symmetrised parts odd under each half, the contact parts even).
\paragraph{What C implies, and what it does not.} C alone allows terms $o\times$(C-odd target structure); these are exactly the $(3\rho)$ finite terms, Eq.~(\dref{eq:sodd}). The absence of $o$ in the logarithms needs, in addition, the factorised form of the logarithms (Section~\dref{sec:proofsketch}): (i) $\ln(\vee/k^+)$: the coefficient is the LO spectrum with $\langle\rho^a\rho^b\rangle$ replaced by its one-step evolution, a property of the projectile alone; a three-charge term can enter it only through a $c$-number invariant tensor $c_{abc}$ (target independent), the only invariant three-tensors are $d^{abc}$ and $f^{abc}$ (null space of the adjoint action, dimension 2 for $SU(3)$, checked), and C-invariance of the projectile evolution allows only the C-even one, $c\propto f^{abc}$, whose contraction with $S^{abc}=d^{abc}g_d$ vanishes. Equivalently: a fully contracted target-independent colour network with one $d$ changes sign under $f\to-f$, hence vanishes. The target factor $\mathcal O_{\rm LO}$ is C-even. (ii) $\ln(k^+/\Lambda)$: the coefficient is $-H_{\rm JIMWLK}\mathcal O_{\rm LO}$ at fixed $n$; the projectile enters only through $n$. So no logarithm can multiply $o$. The explicit cancellation confirms this and is stronger: the termwise $d$-projections of the three-$\rho$ logarithmic terms are C-odd target structures of size $1.6$--$32$ per bracket, which C alone would allow (they come with $o$), and they cancel to $<6\times10^{-15}$ (\texttt{fix\_dabc}); the whole symmetrised tensor cancels, i.e.\ also the C-even $f^{abc}g_f$ part, which C does not forbid. The dipole illustrates (i): the one-step evolution is built from the $q\bar qg$ charge products of Sec.~\ref{r3c:d}, which are symmetric under $q\leftrightarrow\bar q$, so the BFKL term of a dipole and of its C-conjugate are identical while $o$ changes sign.

% ------------------------------------------------------------------------------------------
\subsection{Numerical checks}\label{r3c:num}
Scripts in \texttt{notes/scripts/ref3\_cubic/}, explicit $SU(N_c)$ generators, the 13 tensors transcribed from Appendix~\dref{app:fin_sym} (their $\Delta$ agree with Eq.~\dref{eq:Deltas} and the \texttt{fix\_dabc} table).
\begin{itemize}
\item \texttt{quark\_dipole.py}: states quark, antiquark, dipole, conjugate dipole; $n$, $o$, $g_f$, $S^{abc}$ vs (A2): $\le10^{-15}$; $o_q$ and $f^{abc}t^at^bt^c$ identities exact; exact $\langle C\rho\rho\rho\rangle$ vs symmetrised $+$ contact for all 13 structures and all assignments, $N_c=2,3,4$: largest relative deviation $5.4\times10^{-14}$ ($SU(2)$), $3.0\times10^{-14}$ ($SU(3)$), $2.8\times10^{-16}$ ($SU(4)$) (values that vanish identically, e.g.\ I.II $N_{2b}$ at a triple coincidence, measured against $\max|C|$); C-split (\ref{r3c:Csplit}): $\le8\times10^{-15}$; ordering test as quoted in Sec.~\ref{r3c:b}.
\item \texttt{singlequark\_table.py}: the table of Sec.~\ref{r3c:quark} vs the exact trace, $SU(3)$, $SU(4)$, 4 targets: $\le7\times10^{-13}$ (relative; vanishing entries against $\max|C|$).
\item \texttt{magnitudes.py}: colour weights (table of Sec.~\ref{r3c:dipole}) and weak-target scaling.
\item \texttt{emission.py}, \texttt{cparity.py}, \texttt{ls\_check.py}: Secs.~\ref{r3c:d}, \ref{r3c:e}, \ref{r3c:c}.
\end{itemize}
\paragraph{Not established.} (1) The kernels at the triple coincidence (single quark, local $o$) were not power-counted. (2) No kernel-level comparison with LS/MS (no momentum-space form of the finite terms). (3) The $k$-parity disagreement with LS is not resolved; the prescription of $\mf B_3$ (\texttt{ref3\_prescriptions}) is a candidate, not a demonstrated resolution. (4) The magnitudes are colour weights; the size of the actual contributions depends on the kernels.

% ------------------------------------------------------------------------------------------
\subsection{Drop-in text for the paper}\label{r3c:dropin}
\begingroup\def\v#1{\mathbf{#1}}\providecommand{\Kv}{\mathcal K}\renewcommand{\eqref}[1]{(\texttt{\detokenize{#1}})}\renewcommand{\ref}[1]{\texttt{\detokenize{#1}}}\renewcommand{\cite}[1]{[\texttt{\detokenize{#1}}]}
\noindent Labels and citations are printed as their keys; in the source they are the paper's.

\subsubsection*{(D1) \texttt{sec\_finite.tex}, Sec.~\ref{sec:fin_singlet}: after ``\dots\ the $(3\rho)$ equations are absent and the $(2\rho)$ equations remain.)''}
\noindent\fbox{\parbox{0.97\textwidth}{\small
The split is unique. The order of the charges in each row is not a choice: it is the order in which the coefficients of $\bar\Omega^\dagger\Omega$ and of the conjugate amplitude act, fixed by the light-cone time ordering, Eq.~\eqref{eq:lcpt} (e.g.\ $[\mf A^{(3)}]_{\rho\rho\rho}=\mf A\mf N_2$). The decomposition~\eqref{eq:symmetrise} is an operator identity, and its symmetrised term is the same operator for every ordering of the three charges: the ordering enters only through the contact terms. Under (A1) and (A2) the two parts have an ordering-free meaning: since charge conjugation of the projectile, $\rho^at^a\to-(\rho^at^a)^T$, reverses $o$ and leaves $n$ unchanged, the symmetrised part is $\frac12[\langle\cdot\rangle_P-\langle\cdot\rangle_{P^{\rm C}}]$ and the contact part $\frac12[\langle\cdot\rangle_P+\langle\cdot\rangle_{P^{\rm C}}]$, and by C-invariance these are also the parts odd and even under $V\to V^*$. Section~\ref{sec:check_cubic} shows the split on a single quark and on a dipole, computed exactly in charge space.
}}

\subsubsection*{(D2) \texttt{sec\_checks.tex}: new subsection after Sec.~\ref{sec:check_BFKL}}
\noindent\fbox{\parbox{0.97\textwidth}{\small
\subsection*{Cubic terms on exact charge states: a quark and a dipole}\label{sec:check_cubic}
For a single quark, $\rho^a(\v x)=t^a\delta^{(2)}(\v x-\v x_0)$ in a colour-averaged state, the three-point function is exact: $\langle\rho^a\rho^b\rho^c\rangle=\frac1{N_c}{\rm tr}(t^at^bt^c)=\frac{1}{4N_c}(d^{abc}+if^{abc})$ at $\v x_0$, so that (A2) holds with $o=\frac{(N_c^2-4)(N_c^2-1)}{4N_c^2}=C_F\frac{N_c^2-4}{2N_c}$ ($d^{abc}t^at^bt^c=C_F\frac{N_c^2-4}{2N_c}$) and $n=C_F$. Each cubic colour structure then has the unique value $\frac{1}{4N_c}[\Delta+i\Phi]$ for an ordered product and $\frac{1}{4N_c}\Delta$ for a product with an anticommutator (${\rm tr}\,t^a\{t^b,t^c\}=\frac12d^{abc}$), with $\Delta=\sum C^{abc}d^{abc}$ and $\Phi=C^{abc}f^{abc}$ at $\v x_0$; Eq.~\eqref{eq:cubicavg} reproduces it term by term: the symmetrised part is the $d^{abc}$ half, the contact part the $f^{abc}$ half, and for an antiquark the first changes sign and the second does not. For a colour-singlet dipole, quark at $\v x_0$ and antiquark at $\v x_1$, $n=C_F\,s_is_j$ and $o=o_q\,s_is_js_k$ with $s=\pm1$ at the quark and the antiquark; the symmetrised part of a cubic term carries $s_is_js_k$ and is the part odd under the exchange of quark and antiquark, i.e.\ under the C-conjugation of the dipole, while the contact part carries products of two $s$'s and is even. The $k$-odd cubic term of a dipole is therefore the part of the spectrum that distinguishes the quark from the antiquark end; it needs a C-odd target and vanishes for an orientation-averaged dipole, whereas the $k$-even cubic (contact) part does not. We evaluated all thirteen cubic colour structures with explicit $SU(2)$, $SU(3)$, $SU(4)$ matrices for the quark, the antiquark, the dipole and the conjugate dipole, all assignments of the charge points, and random Wilson lines: the exact ordered expectation values agree with the split~\eqref{eq:cubicavg} to $5\times10^{-14}$. For a quark or a dipole neither part is enhanced: both are of the order of the quadratic terms (colour weights relative to the LO weight, dense random $SU(3)$ target: about $0.12$ for the symmetrised and $0.8$ for the contact part; for a weak target the symmetrised part is suppressed by one power of the target field). The language of a ``first saturation correction'' applies to the symmetrised part only for a coherent dense projectile with large $o$.
}}

\subsubsection*{(D3) \texttt{sec\_checks.tex}, Sec.~\ref{sec:check_BFKL}: the sentence with the emission probability and Eq.~(7.2)}
\noindent Replace ``and the probability to emit the additional gluon at $\v z$ is $\frac{g^2}{4\pi^3}\ln\frac{\vee}{k^+}C_F\,K_D(\v x_1,\v x_2,\v z)$.'' by\par\smallskip
\noindent\fbox{\parbox{0.97\textwidth}{\small
and the probability to emit the additional gluon at $\v z$ is $\frac{g^2C_F}{4\pi^3}\ln\frac{\vee}{k^+}K_D(\v x_1,\v x_2,\v z)=\frac{\alpha_sC_F}{\pi^2}\ln\frac{\vee}{k^+}K_D(\v x_1,\v x_2,\v z)$, $-2\mf N_2$ of Eq.~\eqref{eq:N2} for the dipole. With the LO spectrum of the dipole, $\frac{d^3N_{\rm LO}}{d^3k}=\frac{\alpha_s}{4\pi^4k^+}\mathcal D_{12}$ by Eq.~\eqref{eq:LOsinglet}, real minus virtual gives, exactly at finite $N_c$,
}}\par\smallskip
\noindent and in Eq.~(7.2) replace ``$\times(\text{LO prefactor})$'' by ``$\times\frac{\alpha_s}{4\pi^4k^+}$''.

\subsubsection*{(D4) \texttt{sec\_rapidity.tex}, Sec.~\ref{sec:threerho}: replace ``This is required, since the LO spectrum depends on the projectile only through the C-even two-point function $n$, and its one-step evolution cannot produce the C-odd $o$.''}
\noindent\fbox{\parbox{0.97\textwidth}{\small
That no logarithm can multiply $o$ follows from charge conjugation. C acts as $\rho^at^a\to-(\rho^at^a)^T$, i.e.\ $\rho^a\to-\eta_a\rho^a$ with $(t^a)^T=\eta_at^a$, on the soft gluons as $A^a\to-\eta_aA^a$ and on the target as $V\to V^*$, $U^{ab}\to\eta_a\eta_bU^{ab}$; since $\eta_a\eta_b\eta_cf^{abc}=-f^{abc}$ and $\eta_a\eta_b\eta_cd^{abc}=d^{abc}$, it preserves Eq.~\eqref{eq:rho} and all couplings, and the spectrum is invariant under the joint transformation. Under C, $n$ is even and $o$ odd, $\mf D$, $\mf F$, $\mf Q$ are even and $\mf G$, $\mf Q_d$ odd. By the factorisation of Section~\ref{sec:proofsketch}, the coefficient of $\ln(\vee/k^+)$ is the LO spectrum with the projectile correlator replaced by its one-step evolution, which does not involve the target: a three-charge term can enter it only through a $c$-number invariant tensor, $d^{abc}$ or $f^{abc}$, and C-invariance of the evolution admits only $f^{abc}$, whose contraction with Eq.~\eqref{eq:A2d} vanishes; the target enters through the C-even $\mathcal O_{\rm LO}$. The coefficient of $\ln(k^+/\Lambda)$ is $-H_{\rm JIMWLK}$ acting on the target operator of the LO spectrum, in which the projectile enters only through $n$. Hence neither logarithm can multiply $o$; an $o$-dependence requires a C-odd target structure and is non-logarithmic, Eq.~\eqref{eq:sodd}. The explicit cancellation confirms this: individual terms have $d^{abc}$ projections of order one, with C-odd target factors, which C alone would allow, and they cancel in the sum; the cancellation of the whole symmetrised tensor also removes an $f^{abc}g_f$ part, which is C-even and not forbidden by C.
}}

\subsubsection*{(D5) \texttt{sec\_finite.tex}, after Eq.~\eqref{eq:sodd}: new sentence}
\noindent\fbox{\parbox{0.97\textwidth}{\small
The four-Wilson-line invariants cancel between Rows~II.I and~III.II: both have the kernel $\mathfrak K_{\mathcal T}(w,w')$, Row~III.II with its complex conjugate is $-4\sin(\v k\cdot(\v w-\v w'))\mathfrak K_{\mathcal T}\Delta_{\rm III.II}$, and $\Delta_{\rm II.I}+\Delta_{\rm III.II}=\mf G(x,x',z)+\mf G(x,z,w')$, so that the two lines combine into $-\frac{\mathcal C_{\rm f}}{N_c^2-1}\frac{N_c}{N_c^2-4}\int4\sin(\v k\cdot(\v w-\v w'))\,\mathfrak K_{\mathcal T}(w,w')\,[\mf G(x,x',z)+\mf G(x,z,w')]\,o(x,y,x')$. Likewise the $\mf Q$ of the contact terms of Eqs.~\eqref{eq:sIII1} and~\eqref{eq:sIIIII} cancel. The odd part therefore contains at most three Wilson lines, as the classical results~\cite{McLerran:2016snu,Li:2021ntm}.
}}

\subsubsection*{(D6) \texttt{sec\_checks.tex}, paragraph ``Odd terms'': replace the last sentence (``A comparison of the three-$\rho$ terms \dots\ have not been carried out.'')}
\noindent\fbox{\parbox{0.97\textwidth}{\small
The symmetrised three-$\rho$ terms are the counterpart, for non-commuting charges, of the first saturation correction $M_{(1)}^*M_{(3)}+{\rm c.c.}\propto\rho_P^3$ of Li and Skokov~\cite{Li:2021ntm}, whose charges are classical; the contact terms have no classical counterpart. The structures agree: their two-Wilson-line term has a vanishing $d^{abc}$ projection, as the corresponding structures here, the odd part comes from three Wilson lines and is C-odd in the target, and in our result the four-line invariants cancel between Rows~II.I and III.II. Li and Skokov, however, find also a part of $M_{(1)}^*M_{(3)}$ that is even in $\v k$ configuration by configuration, and its $d^{abc}$ projection is not removed by symmetry, whereas our symmetrised cubic terms are purely odd. A candidate for the difference is the prescription of the vanishing energy denominators of $\mf B_3$, Section~\ref{sec:fin_pv}: the adiabatic out-state dressing adds to Row~I.III a term with a real kernel, whose symmetrised part is even in $\v k$. A comparison of the kernels requires the finite terms in momentum space and has not been carried out.
}}

\subsubsection*{(D7) Reply to the referee (Major 7 and the two error items), for the response letter}
\noindent\fbox{\parbox{0.97\textwidth}{\small
\emph{Ordering.} The order of the charges is fixed by the light-cone time ordering and by the position of the conjugate amplitude; it is not a choice. The split into a symmetrised part and contact terms is an operator identity; the symmetrised part is the same for every ordering, and under (A2) the two parts are the C-odd and C-even parts of the projectile, $\frac12[\sigma(P)\mp\sigma(P^{\rm C})]$, which makes the split unique and physical. \emph{Quark and dipole.} We have added Section~\ref{sec:check_cubic}: for a quark the symmetrised and contact parts are the $d^{abc}$ and $if^{abc}$ halves of $\frac1{N_c}{\rm tr}\,t^at^bt^c$; for a dipole the $k$-odd part is the part odd under the exchange of quark and antiquark; all thirteen structures were verified in exact charge space. \emph{Li--Skokov, McLerran--Skokov.} See the new text in Section~\ref{sec:checks}: agreement on the colour structures (and a simplification of our result: the four-line invariants cancel), a difference in the $\v k$-even part that we trace to the prescription question of Major~5; the kernel comparison requires a momentum-space form and is left open. \emph{Section~7.3.} The manuscript has $\frac{g^2}{4\pi^3}C_FK_D=\frac{\alpha_sC_F}{\pi^2}K_D$; we now define the LO prefactor $\frac{\alpha_s}{4\pi^4k^+}$ of Eq.~(7.2), which was probably the source of the factor $\pi$. \emph{Section~4.3.} Rewritten with C-parity as the reason and the cancellation as confirmation.
}}
\par\endgroup
